\documentclass[11pt]{article}
\usepackage[utf8]{inputenc}
\usepackage{amsmath, amssymb, amsthm}
\usepackage[margin=1in]{geometry}

\theoremstyle{definition}
\newtheorem{definition}{Definition}
\newtheorem{theorem}{Theorem}
\newtheorem{lemma}{Lemma}

\theoremstyle{remark}
\newtheorem{remark}{Remark}

\title{Lecture 14: Mild Solution Analysis and Rough Integral Estimates}
\date{15.01.2026}
\author{Tien Anh Vu}

\begin{document}
\maketitle

Good morning, class. Let us begin today's lecture on Stochastic Partial Differential Equations. We will go through the Lecture 14 notes in detail, focusing on the functional estimates and the rough path integration techniques that arise in our evolution equations.

\section*{1. Mild Decomposition of the Process}
\subsection*{1.1. Decomposition Formula}
We begin by placing the equation in the mild solution framework for our stochastic partial differential equation. The solution process $u(t,x)$ is decomposed into three components:
\[
u(t,x)=v(t,x)+X_t(x)+e^{-t}P_t(u_0-X_0).
\]
This decomposition separates a smoother remainder, a rough noise component, and the damped semigroup evolution of the shifted initial data.

\subsection*{1.2. Mild Form and Operator Splitting}
Next, we write the integral equation for the remainder process $v(t,x)$. It is given by the sum of two integral operators:
\[
v(t,x)=\int_0^t e^{-(t-s)}P_{t-s}\bigl(g(u(s,\cdot))\partial_x v(s,\cdot)+f(u(s,\cdot))\bigr)(x)\,ds
 +\int_0^t\int_0^{2\pi} e^{-(t-s)}P_{t-s}(x-y)g(u(s,y))\,dX_s(y)\,ds.
\]
We label the first time-integral term by $M_{T,X}^{(1)}v(t,x)$ and the second space-time term by $M_{T,X}^{(2)}v(t,x)$. Explicitly,
\[
M_{T,X}^{(1)}v(t,x)=\int_0^t e^{-(t-s)}P_{t-s}\bigl(g(u(s,\cdot))\partial_x v(s,\cdot)+f(u(s,\cdot))\bigr)(x)\,ds,
\]
$M_{T,X}^{(1)}v(t,x)$ is the deterministic contribution: it collects the drift $f(u(s,\cdot))$ and transport term $g(u(s,\cdot))\partial_x v(s,\cdot)$ over past times and propagates them by $e^{-(t-s)}P_{t-s}$.

and
\[
M_{T,X}^{(2)}v(t,x)=\int_0^t\int_0^{2\pi} e^{-(t-s)}P_{t-s}(x-y)g(u(s,y))\,dX_s(y)\,ds.
\]
$M_{T,X}^{(2)}v(t,x)$ is the rough forcing contribution: it integrates $g(u(s,y))$ against the spatial rough signal and propagates the result by the damped kernel $e^{-(t-s)}P_{t-s}(x-y)$.

\section*{2. Two Main Steps of the Analysis}
Our analysis proceeds in two main steps.

\subsection*{2.1. Step 1: The First Estimate}
Our first objective is to establish a rigorous estimate for the operator $M_{T,X}^{(1)}$. We consider the supremum over times $t\leq T$ of the $C^1$ norm of $M_{T,X}^{(1)}v(t)$ and obtain
\[
\sup_{t\le T}\|M_{T,X}^{(1)}v(t)\|_{C^1}\le C_k\bigl(1+\|X\|_\alpha\bigr).
\]
This is a classical evolution estimate: the $C^1$ norm up to time $T$ is controlled by $1+\|X\|_\alpha$.

\subsection*{2.2. Step 2: The Main Ingredient}
In Step 2, we control $M_{T,X}^{(2)}$ in a $C^1$-type norm. Although no explicit derivative appears in its definition, we must estimate its spatial derivative. For $\alpha\in(\frac13,\frac12)$, the key bound is:
\begin{equation}\label{eq:step2-rough-estimate}
\left|\partial_x\int_0^{2\pi} e^{-t}P_t(x-y)Y_y\,dX_s(y)\right|
\le C t^{\alpha/2-1} \|Y\|_{X,\alpha}\, |||\mathbb X|||_{\alpha}.
\end{equation}
This estimate is uniform in $x\in[0,2\pi]$ and isolates the time-singular factor $t^{\alpha/2-1}$.

\section*{3. Context: Classical It\^o Isometry vs Rough Path Integration}
\subsection*{3.1. Classical It\^o Isometry}
We then transition to the stochastic integration properties. More precisely, the left-hand side of \eqref{eq:step2-rough-estimate},
\[
\left|\partial_x\int_0^{2\pi} e^{-t}P_t(x-y)Y_y\,dX_s(y)\right|,
\]
is not itself justified by the It\^o isometry, because it is a spatial rough integral rather than a classical stochastic integral in time. The It\^o isometry is mentioned here only as the classical comparison principle. The foundational equation for this isometry is:
\[
\mathbb E\!\left(\left(\int_0^s dW_r\right)^2\right)=\mathbb E\!\left(\int_0^s dr\right).
\]
This isometry controls classical stochastic integrals. By contrast, \eqref{eq:step2-rough-estimate} is controlled by rough-path estimates.

\subsection*{3.2. Rough Path Definition of the Spatial Integral}
Finally, we must precisely define the spatial rough integral $\int_0^{2\pi} Z_y\,dX(y)$, that is, a rough path integral over space. The integral follows the formal limit:
\[
\int_0^{2\pi} Z_y\,dX(y)
:=
\lim_{|P|\to0}
\sum_{[x,y]\in P}
\bigl(Z_x(X_y-X_x)+Z_x'\,\mathbb X(x,y)\bigr).
\]
The rough integral is defined as the partition limit of the Riemann--Stieltjes term plus the second-order correction term.

Here, $X$ is the rough driving signal, and $Z$ is the function or path that we want to integrate against $X$.

The critical observation concluding these notes is that this formulation does not require $Z_x$ to be differentiable. Therefore, even if the classical derivative of $Z_x$ does not exist, the rough path framework bypasses this limitation and allows us to maintain mathematical rigor throughout the integration process.

\section*{4. Integrand Properties and Rough Bounds}
\subsection*{4.1. The Gubinelli Derivative and the Remainder}
We now focus on the integrand structure used in the rough expansion. Here, $Z_x'$ is the Gubinelli derivative, and the remainder is
\[
R_{x,y}=Z_y-Z_x-Z_x'(X_y-X_x).
\]
$R_{x,y}$ measures the deviation from first-order control by the increment $X_y-X_x$.

As noted in the margin, $Z'$ is not uniquely defined and is not classically differentiable. The framework is valid when $R_{x,y}$ is $2\alpha$-H\"older, namely
\[
\|R\|_{2\alpha}<\infty.
\]
This boundedness gives the higher-order control needed for the rough expansion.

\subsection*{4.2. Contrast with Classical Derivatives}
To draw a sharp contrast with classical calculus, we recall the standard definition of a derivative of one function with respect to another:
\[
\frac{dg}{df}(x)=\lim_{h\to0}\frac{g(x+h)-g(x)}{f(x+h)-f(x)}.
\]
Classically, this derivative is defined by a difference-quotient limit; rough paths replace that requirement by controlled expansions.

\subsection*{4.3. A Fundamental Rough Integral Bound}
This leads us to a fundamental bound for the rough integral, which is written as:
\[
\left\|\int_0^\cdot (Z_y-Z_0)\,dX_y\right\|_{\alpha}
\le C\bigl(\|X\|_{\alpha}\|R\|_{2\alpha}+\|\mathbb X\|_{2\alpha}\|Z'\|_{\alpha}\bigr).
\]
This bound controls the rough integral by the driver regularity and the controlled-path data $(R,Z')$.

\section*{5. Estimates for $M_{T,X}^{(2)}$ and Rough Path 2.0}
\subsection*{5.1. The Heat Semigroup Derivative Estimate}
This subsection refines Step 2 by deriving the explicit spatial-derivative estimate used to bound $M_{T,X}^{(2)}$ in the target norm:
\[
\left|\partial_x\int_0^{2\pi} P_{t-s}(x-y)Y_y\,dX(y)\right|
\le C_K(1+|||X|||)^3\underbrace{\bigl(1+s^{(3-2\alpha)/2}\bigr)(t-s)^{\alpha/2}}_{\sim\,(t-s)^{\alpha/2-1}}.
\]
The key point is the explicit time singularity generated by the derivative of the heat kernel.

\subsection*{5.2. Origin of the $(t-s)$ Singularity}
To understand precisely where this $(t-s)$ singularity originates, we make explicit that the spatial derivative in the $C^1$-estimate falls on the heat kernel. Formally,
\[
\partial_x\int_0^{2\pi} P_{t-s}(x-y)Y_y\,dX(y)
=
\int_0^{2\pi} \partial_x P_{t-s}(x-y)Y_y\,dX(y).
\]
Thus the relevant object is the spatial derivative of the heat kernel itself. Writing the kernel in the form
\[
P_{t-s}(x-y)\sim \frac{1}{\sqrt{t-s}}e^{-\frac{\|x-y\|^2}{2(t-s)}},
\]
the derivation is:
\[
\partial_x\!\left(\frac{1}{\sqrt{t-s}}e^{-\frac{\|x-y\|^2}{2(t-s)}}\right)
=-\frac{x-y}{t-s}\frac{1}{\sqrt{t-s}}e^{-\frac{\|x-y\|^2}{2(t-s)}}.
\]
This computation is the source of the singular $(t-s)$ factor in the derivative estimate.

Putting these elements together, we arrive at the critical estimate for our second integral operator, $M_{T,X}^{(2)}$. The bound is:
\[
\|M_{T,X}^{(2)}v\|_{1,T}
\le C_K(1+|||X|||)^3
\sup_{t\le T}\int_0^t (t-s)^{\alpha/2-1}\,ds
< T^{\alpha/2}.
\]
Hence $\|M_{T,X}^{(2)}v\|_{1,T}$ remains finite because the time singularity is integrable.

\subsection*{5.3. Rough Path 2.0 and Spatial Regularization}
Finally, we transition to the topic I have labeled ``Rough path 2.0'', where we consider a random field with a spatial parameter. We introduce the regularized stochastic partial differential equation model:
\[
du_{\varepsilon}
=\partial_{xx}^2u_{\varepsilon}\,dt+f(u_{\varepsilon})\,dt+g(u_{\varepsilon})\partial_xu_{\varepsilon}+\sigma\circ Q_{\varepsilon}\,dW_t.
\]
This regularized equation combines diffusion, reaction, transport, and a smoothed noise input.

This spatial smoothing operator $Q_{\varepsilon}^u$ acts on a function $\xi$ via convolution. The defining equation is:
\[
Q_{\varepsilon}^u(x)=\int_0^{2\pi}\phi_{\varepsilon}(x-y)u(y)\,dy\to u(x).
\]
$Q_{\varepsilon}^u$ is a spatial mollification of $u$, and the limit $\varepsilon\to0$ recovers the unsmoothed field.


\section*{6. Discretization and Biological Models}
\subsection*{6.1. Discretization and Space-Time White Noise}
We now discuss the numerical issue behind singular SPDEs. Any computation uses a finite discretization in time and space, so the algorithm only sees projected grid data, effectively a smooth approximation of the noise. In particular, a naive scheme does not directly simulate an SPDE driven by ideal space-time white noise; it simulates a regularized model.

This explains why refinements of the mesh may produce unstable fluctuations, and why in some examples the limiting object can depend on the chosen discretization. Rough path theory addresses exactly this dependence: it tracks continuity with respect to the driving signal, identifies diverging correction terms, and provides convergence results as the mesh size goes to zero.

Equivalently, one studies smooth approximations of the noise, indexed by $\varepsilon$, and controls the limit as $\varepsilon\to 0$. The practical point is that discretization is part of the model, and rough path analysis is the tool that makes this dependence mathematically tractable.

\subsection*{6.2. The Nagumo Equation and Nerve Modeling}
We construct a coupled system of partial differential equations to mathematically capture the complex biological process of electrical signal propagation along a nerve axon, incorporating both deterministic biology and random physical fluctuations.

We introduce a model representing a nerve fiber, which the notes depict as a physical domain with an internal ``local reaction.'' The governing evolution equation for this system is the stochastic FitzHugh--Nagumo equation. The primary equation for the state variable $u(t,x)$ is given by:
\[
u_t(t,x)=\partial_{xx}u(t,x)+f(u(t,x))-v(t,x)+\xi(t,x).
\]
This equation combines diffusion, local reaction, recovery, and stochastic forcing in a single nerve-model dynamics.


\subsection*{6.3. The Bistable Reaction Dynamics}
We choose a specific nonlinear cubic form for the local reaction term to create a bistable system with a built-in threshold, mirroring how a biological neuron requires a stimulus of a certain minimum strength before it triggers an action potential.

The local reaction term $f(u)$ is defined as a specific cubic polynomial to generate this threshold behavior. The equation is:
\[
u'=f(u)=u(1-u)(u-a).
\]
The cubic has roots at $0$, $a$, and $1$, with stable states at $0$ and $1$ and a threshold at $a$.


\subsection*{6.4. Traveling Waves}
We analyze the spatial profile over time to understand how localized excitation propagates as a traveling wave.

When the system is pushed past the unstable threshold $a$, the dynamics drive the state locally towards the ``white site'' at $1$. The notes illustrate this with a graph of ``traveling wave'' solutions. The diagram displays multiple wave fronts shifting continuously to the right along the $x$-axis, with arrows indicating the state pushing upward to $1$ at the front and decaying backward. The text clarifies that these waves propagate at a ``deterministically constant speed'' across the domain.


\subsection*{6.5. The Recovery Variable and Bump Solutions}
We introduce a secondary, slower differential equation to the system to ensure that after the spatial region reaches its excited state, it is forced to eventually return to its resting state, thereby creating a localized and transient ``bump'' of activity rather than a perpetually excited domain.

To manage these ``isolated regions'' of excitation and prevent the wave from keeping the entire domain at state $1$, the notes indicate we ``need 2nd variable to push it down.'' This is the role of the recovery variable $v(t,x)$. Its evolution is governed by the equation:
\[
\partial_t v(t,x)=\varepsilon\bigl(u(t,x)-\gamma v(t,x)\bigr).
\]
Because $v$ enters the $u$-equation with negative sign, it pulls the state back after excitation and supports localized bump-like pulses.


\section*{7. Observation, Estimation, and Bayesian Updating}
\subsection*{7.1. The Observation Problem and State Space Formulation}
We construct a combined mathematical state vector to capture the complete dynamics of the system because observing only a single variable of a coupled biological or physical system leads to missing information, which destroys the Markov property and prevents us from writing down a closed-form differential equation.

We begin by considering a physical observation process. As noted at the top of the page, in practice, a technique like ``microscopy measures $u$.'' There is also an additive component ``$+g$'' which is tentatively labeled with an arrow as ``measurement?'' suggesting a measurement operator or noise associated with the observation. However, the underlying system requires us to reconstruct both the observed state $u$ and the unobserved, hidden state $v$.

To handle this, we formally define our complete state vector as
\[
X=[u,v].
\]
The state vector is $X=[u,v]$. If $v$ is ignored, the observed component $u$ is not Markovian and does not yield a closed evolution equation.


\subsection*{7.2. The Frequentist Approach to Estimation}
We formulate a direct mathematical inverse problem from our observations to attempt to retrieve the exact underlying state by treating the governing equations as a deterministic mapping, assuming any measurement noise can simply be averaged out over multiple samples.

We formally state our objective:
\[
\text{Est } X \text{ given } Y=G(X,e).
\]
We estimate the hidden state $X$ from measurements $Y=G(X,e)$, where $G$ is the forward model.

The notes first outline the frequentist view: if $G$ is invertible, recover $X$ directly from $Y$.

For a scenario with repeated, additive scalar measurements, we write
\[
Y_n=X+e_n,\qquad n=1,\dots,N.
\]
For repeated scalar observations, $Y_n=X+e_n$ and the estimator is the sample mean:
\[
\hat Y=\frac1N\sum_{n=1}^N Y_n.
\]
This averages independent measurement noise across samples.


\subsection*{7.3. The Bayesian Approach and Likelihood Modeling}
We shift from a direct inversion to a probabilistic framework incorporating prior knowledge and observation noise models because complex inverse problems are often ill-posed or non-invertible, and representing our knowledge as probability distributions allows for a mathematically rigorous quantification of uncertainty.

When $G$ is not easily invertible, we transition to the ``Bayesian approach,'' which the notes break down into two distinct steps.

Step 1 provides a prior model for $X$, denoted $P_X$.

Step ``2) dist of measurement error.'' We must quantify the probability of observing a specific measurement given a known state. This is governed by the conditional probability equation:
\[
P(Y\in dy\mid X=x)=e^{-\ell(x,y)}\,dy.
\]
This likelihood model encodes observation noise via the log-likelihood $\ell(x,y)$.


\subsection*{7.4. The Posterior Distribution}
We combine prior and data through Bayes' theorem to obtain the posterior distribution.

By combining the prior and the likelihood, the notes state this ``yields posterior.'' The fundamental equation for the posterior measure is
\[
P(X\in dx\mid Y=y)=\frac{e^{-\ell(x,y)}P_X(dx)}{\int e^{-\ell(x,y)}P_X(dx)}.
\]
The posterior is the likelihood-weighted prior, normalized by the evidence term.


\subsection*{7.5. Graphical Intuition of the Bayesian Update}
We visually represent the multiplication of probabilities to build a concrete geometric understanding of how a sharp, possibly multi-modal prior belief is modulated and shifted by the broader distribution of incoming measurement data.

At the bottom of the page, I have sketched three graphs to illustrate this multiplicative update.

The top-left graph is labeled ``prior.'' The horizontal axis represents the state space, and the vertical axis represents probability density. The sketch shows two distinct, sharp peaks, indicating a strong prior belief that the state exists in one of two specific, narrow configurations.

There is a plus sign leading to the top-right graph, labeled ``measurement density.'' This graph displays a single, smooth, broad curve---similar to a wide Gaussian distribution---representing the information provided by the noisy observation $Y$.

The bottom graph illustrates the resulting posterior distribution, which is the product of the prior and the measurement density, subsequently normalized. The sketch demonstrates that the posterior retains the sharp, multi-modal characteristics of the prior peaks, but their relative heights and exact shapes are now modulated and scaled by the broad envelope of the measurement density curve.


\section*{8. Continuous-Time Filtering and the Kallianpur--Striebel Formula}
Let us continue our lecture series. We now turn to the fifth page of the Lecture 14 notes, where the continuous-time filtering problem is formulated and the Kallianpur--Striebel formula is introduced.

\subsection*{8.1. The Signal and Measurement Equations}
We begin by establishing the governing stochastic differential equations for the system. The first equation defines the hidden signal, marked $(S)$ in the notes:
\[
dX_t = B(X_t)\,dt + C(X_t)\,dW_t.
\]
Expressed in words, the differential of the state vector $X_t$ is equal to the drift vector field $B(X_t)$ multiplied by $dt$, plus the diffusion matrix $C(X_t)$ integrated against the Wiener process $dW_t$. The state $X_t$ evolves in $\mathbb R^d$. As noted in the margin, one assumes global Lipschitz conditions on the coefficients to guarantee existence and uniqueness of the solution.

The second equation governs the observation, labeled $(O)$:
\[
dY_t = H(X_t)\,dt + R\,dV_t,
\qquad
Y_0=0.
\]
This says that the differential of the observation vector $Y_t$ is equal to the observation function $H(X_t)$ multiplied by $dt$, plus the noise coefficient matrix $R$ integrated against a separate Wiener process $dV_t$. The observation process takes values in $\mathbb R^p$.

The notes stress two structural features of the problem. First, the true state $X$ is not observed directly, so its initial condition is not available to us. Second, the observation dimension satisfies $p\ll d$, which reflects the fact that the data are much lower-dimensional than the hidden state.

Finally, we formalize the history of the observations by
\[
\mathcal Y_{0:t}=\sigma(Y_s:s\le t).
\]
In words, $\mathcal Y_{0:t}$ is the sigma-algebra generated by the observation process up to time $t$, i.e.\ the information available up to time $t$.

\subsection*{8.2. Formulation of the Filtering Problem}
With the system defined, we state the filtering problem:
\[
\eta_t(A)=\mathbb P(X_t\in A\mid \mathcal Y_{0:t}),
\qquad
A\in\mathcal B(\mathbb R^d).
\]
In words, $\eta_t(A)$ is the conditional probability that the hidden state $X_t$ lies in the set $A$, given the observation history up to time $t$.

\subsection*{8.3. The Kallianpur--Striebel Formula}
To solve this problem, we introduce the central theorem of this section. Let $(\bar X_t)$ be a copy of $(X_t)$ such that $\bar X$ is independent of the pair of processes $(X,V)$.

\begin{theorem}[Kallianpur--Striebel]
For every $A\in\mathcal B(\mathbb R^d)$,
\[
\eta_t(A)
:=
\mathbb E\bigl(\mathbf 1_A(X_t)\mid \mathcal Y_{0:t}\bigr)
=
\frac{\mathbb E_{\bar X}\bigl(\mathbf 1_A(\bar X_t)Z_t(\bar X,Y)\bigr)}
{\mathbb E_{\bar X}\bigl(Z_t(\bar X,Y)\bigr)},
\]
where
\[
Z_t(\bar X,Y)
=
\exp\left(
\int_0^t \left\langle R^{-1}H(\bar X_s),dY_s\right\rangle
-\frac12\int_0^t \left\langle R^{-1}H(\bar X_s),H(\bar X_s)\right\rangle ds
\right).
\]
\end{theorem}

This formula rewrites the conditional law of the hidden state as a ratio of expectations over the independent copy $\bar X$, with the observed path $Y$ held fixed.
Here $Z_t(\bar X,Y)$ is the likelihood weight attached to the candidate path $\bar X$ given the observed data $Y$ up to time $t$.

\subsection*{8.4. Setting Up the Proof}
At the bottom of the page, the proof begins by fixing a realization of the signal path $X(\omega)$. Once the signal path is frozen, the observation equation becomes a deterministic signal plus additive Gaussian noise. This is the point at which one can derive the conditional distribution of the observations given the state and apply change-of-measure arguments such as Girsanov's theorem.

\section*{9. Likelihood Densities, Measure Change, and the Zakai Equation}
Let us resume our lecture. We now turn to the sixth page of the notes, which concludes the derivation of the continuous-time filtering problem and connects the probabilistic formulation to stochastic partial differential equations.

\subsection*{9.1. The Likelihood Density and the Radon--Nikodym Derivative}
We continue the proof by computing the conditional law of the observations when the signal path is held fixed. In the notes, $X$ is marked as ``a realization,'' meaning that for this step we regard the hidden path as given. The relevant object is the Radon--Nikodym derivative of the conditional observation measure with respect to the pure noise measure. The formula is written in the notes as
\[
\frac{dP_{Y\mid X}}{dP_V}
=
\exp\left(
\int_0^t \left\langle R^{-1}H(\bar X_s),dV_s\right\rangle
-\frac12\int_0^t \left\langle R^{-1}H(\bar X_s),H(\bar X_s)\right\rangle ds
\right).
\]
In words, the Radon--Nikodym derivative of the conditional law of the observations with respect to the Wiener measure $P_V$ is given by an exponential weight involving a stochastic integral and its compensating quadratic term.

The notes also stress an important distinction from the finite-dimensional Bayesian formula. One may write heuristically $e^{-\ell(x,y)}\,dy$, but here the correct object is not a Lebesgue density. It is a Radon--Nikodym density with respect to the Wiener measure $P_V$. In path space, the natural reference measure is a Gaussian process measure, not Lebesgue measure.

\subsection*{9.2. Independence and Measure Transformation}
To justify the separation of state and observation noise, the notes invoke Girsanov's theorem. Under the corresponding reference measure, the observation drift is removed and the observation process behaves like pure noise, independent of the hidden signal. This yields the product structure
\[
P_{(X,V)}=P_X\otimes P_V.
\]
This product form justifies integrating over the independent state copy while treating the observation path as fixed.

The notes also remark that the full proof should be checked in the script, since the complete measure-theoretic argument is longer than what fits into the handwritten summary.

\subsection*{9.3. The Unnormalized Measure and the SPDE Connection}
The lower half of the page asks: what is the relevance to SPDEs? The bridge is the unnormalized conditional measure, denoted $\tilde \eta_t$. For a test function $f$, it is defined by
\[
\tilde \eta_t(f)
=
\int f\,d\tilde \eta_t
=
\mathbb E_{\bar X}\bigl[f(\bar X_t)Z_t(\bar X,Y)\bigr].
\]
$\tilde \eta_t(f)$ denotes the action of the unnormalized filtering measure on the test function $f$.
$\tilde \eta_t(f)$ is exactly the numerator of the Kallianpur--Striebel formula.

\subsection*{9.4. The Zakai Equation and Stochastic Calculus}
To derive the dynamics of $\tilde \eta_t(f)$, one differentiates the product $f(\bar X_t)Z_t(\bar X,Y)$. The notes emphasize that this requires semimartingale theory, since both the state process and the likelihood process are continuous semimartingales. Ordinary calculus does not apply, so one must use the It\^o product rule.

Applying the It\^o product rule to $df(\bar X_t)$ and $dZ_t$, and then taking the expectation $\mathbb E_{\bar X}$, leads to a linear stochastic partial differential equation for the unnormalized filter. As indicated in the notes, this SPDE is the Zakai equation.

Once $\tilde \eta_t$ is known, the normalized filter is recovered by dividing by the total mass:
\[
\eta_t(f)=\frac{\tilde \eta_t(f)}{\tilde \eta_t(1)}.
\]
In words, the true conditional estimate is obtained by normalizing the unnormalized measure at time $t$.

\appendix
\section*{Appendix A}
\subsection*{A.1. Difference Between $P_{t-s}$ and $P_t(x-y)$}
The notation $P_{t-s}$ refers to the heat semigroup as an operator acting on a function, while $P_t(x-y)$ refers to the corresponding heat kernel evaluated at the spatial points $x$ and $y$. These are two ways of describing the same object. More precisely, if $\phi$ is a test function, then
\[
P_{t-s}\phi(x)=\int P_{t-s}(x-y)\phi(y)\,dy.
\]
Thus, $P_{t-s}$ emphasizes the operator that evolves data from time $s$ to time $t$, whereas $P_{t-s}(x-y)$ emphasizes the kernel representation of that operator in space.

\subsection*{A.2. Breakdown of $M_{T,X}^{(1)}v(t,x)$}
Recall that
\[
M_{T,X}^{(1)}v(t,x)=\int_0^t e^{-(t-s)}P_{t-s}\bigl(g(u(s,\cdot))\partial_x v(s,\cdot)+f(u(s,\cdot))\bigr)(x)\,ds.
\]
Its pieces can be displayed as
\[
\begin{aligned}
s &\in [0,t] && \text{past time variable},\\
e^{-(t-s)}P_{t-s} &&& \text{damped heat semigroup from time } s \text{ to time } t,\\
g(u(s,\cdot))\partial_x v(s,\cdot) &&& \text{transport term},\\
f(u(s,\cdot)) &&& \text{drift term},\\
(\cdot)(x) &&& \text{evaluation at the spatial point } x.
\end{aligned}
\]
The mild form comes from Duhamel's principle:
\[
\text{solution at time } t
=
\text{linear evolution of initial data}
+
\int_0^t
\text{semigroup from } s \text{ to } t
\times
\text{forcing at time } s
\,ds.
\]
In the present setting, the forcing is
\[
g(u(s,\cdot))\partial_x v(s,\cdot)+f(u(s,\cdot)),
\]
so the deterministic contribution is exactly
\[
M_{T,X}^{(1)}v(t,x)
=
\int_0^t
e^{-(t-s)}P_{t-s}
\bigl(g(u(s,\cdot))\partial_x v(s,\cdot)+f(u(s,\cdot))\bigr)(x)\,ds.
\]

\subsection*{A.3. Breakdown of $M_{T,X}^{(2)}v(t,x)$}
Recall that
\[
M_{T,X}^{(2)}v(t,x)=\int_0^t\int_0^{2\pi} e^{-(t-s)}P_{t-s}(x-y)g(u(s,y))\,dX_s(y)\,ds.
\]
Its pieces can be displayed as
\[
\begin{aligned}
s &\in [0,t] && \text{past time variable},\\
y &\in [0,2\pi] && \text{spatial integration variable},\\
e^{-(t-s)}P_{t-s}(x-y) &&& \text{damped heat semigroup kernel from } (s,y) \text{ to } (t,x),\\
g(u(s,y)) &&& \text{coefficient multiplying the rough forcing},\\
dX_s(y) &&& \text{integration against the spatial rough signal at time } s,\\
\int_0^{2\pi}\cdots dX_s(y) &&& \text{rough spatial integral at fixed time } s,\\
\int_0^t \cdots ds &&& \text{accumulation over all past times } s.
\end{aligned}
\]
The mild form again comes from Duhamel's principle, but now the forcing term is rough in space. At each fixed time $s$, one first forms the rough spatial integral
\[
\int_0^{2\pi} P_{t-s}(x-y)g(u(s,y))\,dX_s(y),
\]
and then propagates that contribution to time $t$ with the damping factor $e^{-(t-s)}$. Hence the rough contribution is
\[
M_{T,X}^{(2)}v(t,x)
=
\int_0^t
\int_0^{2\pi}
e^{-(t-s)}P_{t-s}(x-y)g(u(s,y))\,dX_s(y)\,ds.
\]

\subsection*{A.4. Meaning of $|||Y|||_{\mathbb X,\alpha}$}
The quantity $|||Y|||_{\mathbb X,\alpha}$ denotes the controlled rough path norm of $Y$ relative to the reference rough path $\mathbb X$ with H\"older exponent $\alpha$. In rough path theory, this means that $Y$ is not viewed in isolation: it is viewed together with a description of how its increments are controlled by the increments of the rough path $\mathbb X$. Thus, $|||Y|||_{\mathbb X,\alpha}$ measures both the H\"older size of $Y$ and its compatibility with the rough structure generated by $\mathbb X$.

\subsubsection*{A.4.1. Controlled Rough Path Norm}
More concretely, if $Y$ is controlled by $X$ with Gubinelli derivative $Y'$, so that the remainder
\[
R^Y(x,y):=Y(y)-Y(x)-Y'(x)\bigl(X(y)-X(x)\bigr)
\]
has better regularity, then a typical controlled rough path norm is
\[
|||Y|||_{\mathbb X,\alpha}
:=
\|Y\|_{C^\alpha}
+
\|Y'\|_{C^\alpha}
+
\|R^Y\|_{C^{2\alpha}}.
\]
Here,
\[
\begin{aligned}
\|Y\|_{C^\alpha}
&:=
\sup_{x\neq y}\frac{|Y(y)-Y(x)|}{|y-x|^\alpha},\\
\|Y'\|_{C^\alpha}
&:=
\sup_{x\neq y}\frac{|Y'(y)-Y'(x)|}{|y-x|^\alpha},\\
\|R^Y\|_{C^{2\alpha}}
&:=
\sup_{x\neq y}\frac{|R^Y(x,y)|}{|y-x|^{2\alpha}}.
\end{aligned}
\]
This formula records precisely that the increment of $Y$ is approximated to first order by $Y'(x)(X(y)-X(x))$, with a remainder that is of higher order.

\subsubsection*{A.4.2. Rough Path Norm of $\mathbb X$}
The rough path norm of $\mathbb X$ itself is typically written as
\[
|||\mathbb X|||_{\alpha}
:=
\|X\|_{C^\alpha}
+
\|\mathbb X\|_{C^{2\alpha}},
\]
where
\[
\|\mathbb X\|_{C^{2\alpha}}
:=
\sup_{x\neq y}\frac{|\mathbb X(x,y)|}{|y-x|^{2\alpha}}.
\]
Thus, $|||\mathbb X|||_{\alpha}$ measures both the $\alpha$-H\"older size of the first-level path $X$ and the $2\alpha$-H\"older size of its second-order component $\mathbb X$.

\subsection*{A.5. H\"older Seminorm and Norm}
Strictly speaking, the quantity
\[
[Y]_{C^\alpha}
:=
\sup_{x\neq y}\frac{|Y(y)-Y(x)|}{|y-x|^\alpha}
\]
is a H\"older seminorm rather than a true norm, because it vanishes on every constant function. A genuine H\"older norm is obtained by adding the uniform size:
\[
\|Y\|_{C^\alpha}
:=
\|Y\|_{L^\infty}
+
[Y]_{C^\alpha}.
\]
Here,
\[
\|Y\|_{L^\infty}
:=
\sup_x |Y(x)|
\]
is the supremum norm, which measures the maximum absolute size of the function $Y$ over its domain.

\subsection*{A.6. Spatial Variable Instead of Time}
In classical stochastic integration, the irregular driving signal is usually indexed by time, for example in an integral of the form
\[
\int_0^t H_s\,dW_s.
\]
Here, the integration variable is the time variable $s$. In the present rough path setting, however, the irregular driving signal is indexed by space, as in
\[
\int_0^{2\pi} Z_y\,dX(y).
\]
Here, the integration variable is the spatial variable $y$. Thus, when we say that we are replacing the temporal variable by a spatial variable, we mean that the rough path integration is performed along space rather than along time.

\subsection*{A.7. Derivation of the Bound in 4.3}
Let
\[
I_t:=\int_0^t (Z_y-Z_0)\,dX_y.
\]
To estimate the $\alpha$-H\"older norm of $I$, we look at an increment over an interval $[s,t]$:
\[
I_t-I_s
=
\int_s^t (Z_y-Z_s)\,dX_y
+
(Z_s-Z_0)(X_t-X_s).
\]
Now assume that $Z$ is controlled by $X$, so that
\[
Z_y-Z_x
=
Z_x'(X_y-X_x)+R_{x,y}.
\]
Applying the rough integral expansion on $[s,t]$ gives
\[
\int_s^t (Z_y-Z_s)\,dX_y
=
Z_s'\,\mathbb X(s,t)+\mathrm{Error}_{s,t},
\]
where the error is controlled by the remainder term and the H\"older regularity of $Z'$. More precisely,
\[
|\mathrm{Error}_{s,t}|
\le
C\Bigl(\|X\|_{\alpha}\|R\|_{2\alpha}
+
\|\mathbb X\|_{2\alpha}\|Z'\|_{\alpha}\Bigr)|t-s|^{3\alpha}.
\]
Since
\[
|(Z_s-Z_0)(X_t-X_s)|
\le
\|Z\|_{C^\alpha}\|X\|_{C^\alpha}|t-s|^\alpha,
\]
the full increment of $I$ is therefore controlled by the same rough path data. Dividing by $|t-s|^\alpha$ and taking the supremum over $s\neq t$ yields
\[
\left\|\int_0^\cdot (Z_y-Z_0)\,dX_y\right\|_{\alpha}
\le
C\Bigl(\|X\|_{\alpha}\|R\|_{2\alpha}
+
\|\mathbb X\|_{2\alpha}\|Z'\|_{\alpha}\Bigr).
\]
Thus, the first term on the right-hand side comes from the remainder $R_{x,y}$, while the second term comes from the second-order component $\mathbb X$ together with the Gubinelli derivative $Z'$.

\subsection*{A.8. Why the Expansion on $[s,t]$ Has the Form $Z_s'\mathbb X(s,t)+\mathrm{Error}_{s,t}$}
If $Z$ is controlled by $X$, then for $y\in[s,t]$ one has
\[
Z_y-Z_s
=
Z_s'(X_y-X_s)+R_{s,y}.
\]
Integrating this identity against $dX_y$ over the interval $[s,t]$ gives
\[
\int_s^t (Z_y-Z_s)\,dX_y
=
\int_s^t Z_s'(X_y-X_s)\,dX_y
+
\int_s^t R_{s,y}\,dX_y.
\]
Since $Z_s'$ is frozen at the left endpoint $s$, it can be pulled out of the first integral, and the remaining term is exactly the second-order rough path increment:
\[
\int_s^t Z_s'(X_y-X_s)\,dX_y
=
Z_s'\,\mathbb X(s,t).
\]
The second integral,
\[
\int_s^t R_{s,y}\,dX_y,
\]
is of higher order and is therefore absorbed into the error term $\mathrm{Error}_{s,t}$. This is why the rough integral expansion takes the form
\[
\int_s^t (Z_y-Z_s)\,dX_y
=
Z_s'\,\mathbb X(s,t)+\mathrm{Error}_{s,t}.
\]

\subsection*{A.9. Derivation of the Error Bound}
Starting again from the controlled expansion
\[
Z_y-Z_s
=
Z_s'(X_y-X_s)+R_{s,y},
\]
we integrate against $dX_y$ on the interval $[s,t]$:
\[
\int_s^t (Z_y-Z_s)\,dX_y
=
\int_s^t Z_s'(X_y-X_s)\,dX_y
+
\int_s^t R_{s,y}\,dX_y.
\]
The first term is the leading rough path contribution,
\[
\int_s^t Z_s'(X_y-X_s)\,dX_y
=
Z_s'\,\mathbb X(s,t).
\]
Hence the error term is what remains after subtracting this leading part. The rough sewing estimate shows that
\[
\left|
\int_s^t (Z_y-Z_s)\,dX_y - Z_s'\mathbb X(s,t)
\right|
\le
C\Bigl(
\sup_{y\in[s,t]}|R_{s,y}|\,|X_{s,t}|
+
\sup_{y\in[s,t]}|Z_y'-Z_s'|\,|\mathbb X(s,t)|
\Bigr).
\]
Now use the H\"older bounds
\[
|R_{s,y}|
\le
\|R\|_{2\alpha}|y-s|^{2\alpha},
\qquad
|X_{s,t}|
\le
\|X\|_{\alpha}|t-s|^{\alpha},
\]
and
\[
|Z_y'-Z_s'|
\le
\|Z'\|_{\alpha}|y-s|^{\alpha},
\qquad
|\mathbb X(s,t)|
\le
\|\mathbb X\|_{2\alpha}|t-s|^{2\alpha}.
\]
Since $|y-s|\le |t-s|$ for $y\in[s,t]$, it follows that
\[
\sup_{y\in[s,t]}|R_{s,y}|\,|X_{s,t}|
\le
\|X\|_{\alpha}\|R\|_{2\alpha}|t-s|^{3\alpha},
\]
and
\[
\sup_{y\in[s,t]}|Z_y'-Z_s'|\,|\mathbb X(s,t)|
\le
\|\mathbb X\|_{2\alpha}\|Z'\|_{\alpha}|t-s|^{3\alpha}.
\]
Combining these two bounds yields
\[
|\mathrm{Error}_{s,t}|
\le
C\Bigl(
\|X\|_{\alpha}\|R\|_{2\alpha}
+
\|\mathbb X\|_{2\alpha}\|Z'\|_{\alpha}
\Bigr)|t-s|^{3\alpha}.
\]
Thus, one part of the error comes from the remainder term $R_{s,y}$, while the other comes from the H\"older regularity of the Gubinelli derivative $Z'$.

\subsection*{A.10. Girsanov's Theorem}
\begin{theorem}[Girsanov]
Let $(\Omega,\mathcal F,(\mathcal F_t)_{t\ge 0},P)$ be a filtered probability space, let $(W_t)_{t\ge 0}$ be an $(\mathcal F_t)$-Brownian motion in $\mathbb R^m$, and let $\theta=(\theta_t)_{t\ge 0}$ be an $\mathbb R^m$-valued progressively measurable process such that
\[
\mathbb E\left[\exp\left(\frac12\int_0^T |\theta_s|^2\,ds\right)\right]<\infty
\qquad \text{for every } T>0.
\]
Define
\[
Z_t
=
\exp\left(
-\int_0^t \langle \theta_s,dW_s\rangle
-\frac12\int_0^t |\theta_s|^2\,ds
\right).
\]
Then $(Z_t)_{t\ge 0}$ is a martingale. If the probability measure $Q$ is defined by
\[
\frac{dQ}{dP}\Big|_{\mathcal F_t}=Z_t,
\]
then the process
\[
\widetilde W_t
=
W_t+\int_0^t \theta_s\,ds
\]
is an $(\mathcal F_t)$-Brownian motion under $Q$.
\end{theorem}

This theorem says that a drift term can be transferred from the stochastic differential equation into the density of a new probability measure. Under the new measure $Q$, the process with drift becomes a Brownian motion again. In the filtering problem, this is the step that removes the observation drift and allows the likelihood to be written as a Radon--Nikodym density with respect to Wiener measure.

\end{document}
